\section{正则映射的连续性}\label{sec:continuity-of-regular-maps}

我们已经引入正则映射, 作为拟射影簇情形中的\textquotedbl{}保结构\textquotedbl{}映射. 既然每个拟射影簇都配有 Zariski 拓扑, 自然会问: 正则映射是否保持这一拓扑的\textquotedbl{}结构\textquotedbl{}? 本节致力于回答这一问题. 我们引入拓扑学中的保结构映射, 即连续映射, 然后证明拟射影簇之间的正则映射关于 Zariski 拓扑是连续的.

先回顾熟悉的连续性概念. 初次接触连续性通常是在微积分中研究函数 $f:\mathbb{R}\to\mathbb{R}$ 时. 在那里, 可以先把连续函数介绍为不用提起笔就能画出图像的函数, 但这种说法相当不正式. 为使之严格, 学生接着要理解连续性的 $\epsilon$-$\delta$ 定义: 对每个 $x_{0}\in\mathbb{R}$ 和每个 $\epsilon>0$, 存在 $\delta>0$, 使 $|x-x_{0}|<\delta$ 蕴含 $|f(x)-f(x_{0})|<\epsilon$.

虽然久而久之 $\epsilon$-$\delta$ 定义会变得直观, 它仍有些繁琐. 幸运的是, 连续性有一个优雅的拓扑刻画: 函数 $f:\mathbb{R}\to\mathbb{R}$ 连续 (按 $\epsilon$-$\delta$ 定义), 当且仅当对每个 Euclid 开集 $V\subseteq\mathbb{R}$, $f^{-1}(V)$ 在 $\mathbb{R}$ 中是 Euclid 开集. 这一刻画的重要性在于, 它引出了连续性的纯拓扑表述. 即使没有度量距离的方法, 这个表述仍有意义, 从而得到以下一般的连续性概念.

\begin{Definition}{连续映射}\leavevmode\label{def:continuous-map-of-topological-spaces}
令 $X,Y$ 为拓扑空间. 若对每个开集 $V\subseteq Y$, 原像 $F^{-1}(V)$ 在 $X$ 中开, 则称映射 $F:X\to Y$ 为\emph{连续的}.
\end{Definition}

虽然通常用开集定义连续性, 正如这里所做的那样, 也可以用闭集重新表述连续性. 特别地 (Exercise~\ref{ex:equivalent-open-and-closed-characterizations-of-continuity}), 映射 $F:X\to Y$ 连续, 当且仅当对每个闭集 $Z\subseteq Y$, $F^{-1}(Z)$ 在 $X$ 中闭. 现在用第二种表述陈述本节的主要结果.

\begin{Theorem}{正则映射是连续的}\leavevmode\label{thm:regular-maps-are-zariski-continuous}
令 $F:X\to Y$ 为拟射影簇之间的正则映射. 若 $Z$ 在 $Y$ 中闭, 则 $F^{-1}(Z)$ 在 $X$ 中闭.
\end{Theorem}
\begin{Proof}
令 $X\subseteq\mathbb{P}^{m}$ 和 $Y\subseteq\mathbb{P}^{n}$ 为拟射影簇, $F:X\to Y$ 为正则映射. 给定闭子簇 $Z\subseteq Y$, 我们要证明 $F^{-1}(Z)$ 在 $X$ 中闭. 由 Lemma~\ref{lem:closedness-is-local}, 只须证明 $X$ 的每个点 $p$ 都包含在某个邻域 $U_{p}$ 中, 使 $F^{-1}(Z)\cap U_{p}$ 在 $U_{p}$ 中闭.

考察点 $p\in X$. 由正则映射的定义, 选取同次数的齐次多项式 $f_{0},\ldots,f_{n}\in K[x_{0},\ldots,x_{m}]$, 使 $p\notin\mathcal{V}(f_{0},\ldots,f_{n})$, 且
\[
  F(a)=[f_{0}(a):\cdots:f_{n}(a)]
  \qquad\text{对所有 }a\in X\setminus\mathcal{V}(f_{0},\ldots,f_{n}).
\]
置 $U_{p}=X\setminus\mathcal{V}(f_{0},\ldots,f_{n})$. 则 $U_{p}$ 在 $X$ 中开且包含 $p$. 考察由 $F_{p}=[f_{0}:\cdots:f_{n}]$ 整体定义的正则函数 $F_{p}:U_{p}\to\mathbb{P}^{n}$. 注意
\[
  F^{-1}(Z)\cap U_{p}=F_{p}^{-1}(Z)=F_{p}^{-1}(\overline{Z}),
\]
其中 $\overline{Z}$ 是 $Z$ 在 $\mathbb{P}^{n}$ 中的射影闭包. 最后一个等式成立, 是因为 $Z$ 在 $Y$ 中闭, 因而 $Z=\overline{Z}\cap Y$. 由于 $\overline{Z}$ 在 $\mathbb{P}^{n}$ 中闭, 存在齐次多项式 $g_{1},\ldots,g_{k}\in K[y_{1},\ldots,y_{n}]$, 使 $\overline{Z}=\mathcal{V}(g_{1},\ldots,g_{k})$. 于是, 对每个 $i=1,\ldots,k$, $g_{i}(f_{0},\ldots,f_{n})$ 都是 $K[x_{1},\ldots,x_{m}]$ 中的齐次多项式. 稍加思考 (Exercise~\ref{ex:preimage-of-projective-vanishing-set-under-polynomial-map}) 可知
\[
  F_{p}^{-1}(\overline{Z})
  =U_{p}\cap\mathcal{V}\bigl(g_{1}(f_{0},\ldots,f_{n}),\ldots,g_{k}(f_{0},\ldots,f_{n})\bigr).
\]
所以 $F^{-1}(Z)\cap U_{p}=F_{p}^{-1}(\overline{Z})$ 在 $U_{p}$ 中闭, 如所需.
\end{Proof}

连续映射是拓扑学中的保结构映射. 因此, 有连续逆映射的连续映射给出了拓扑空间之间的一种自然等价概念, 称为同胚.

\begin{Definition}{同胚}\leavevmode\label{def:homeomorphism-of-topological-spaces}
令 $X,Y$ 为拓扑空间. 从 $X$ 到 $Y$ 的\emph{同胚}是具有连续逆映射 $G:Y\to X$ 的连续映射 $F:X\to Y$. 若 $X,Y$ 之间存在一个同胚, 就称 $X$ 与 $Y$ \emph{同胚}.
\end{Definition}

由于拟射影簇的同构是具有正则逆映射的正则映射, Theorem~\ref{thm:regular-maps-are-zariski-continuous} 蕴含, 拟射影簇的同构 $F:X\to Y$ 是关于 $X,Y$ 上 Zariski 拓扑的同胚. 这一观察体现在下述结果的前两条 (等价的) 断言中.

\begin{Corollary}{同构使子簇一一对应}\leavevmode\label{cor:quasiprojective-isomorphisms-identify-subvarieties}
若 $F:X\to Y$ 为拟射影簇的同构, 且 $Z\subseteq X$, 则下列断言成立.
\begin{enumerate}[label=\arabic*.]
  \item $Z$ 在 $X$ 中开, 当且仅当 $F(Z)$ 在 $Y$ 中开.
  \item $Z$ 在 $X$ 中闭, 当且仅当 $F(Z)$ 在 $Y$ 中闭.
  \item $Z$ 是 $X$ 的子簇, 当且仅当 $F(Z)$ 是 $Y$ 的子簇.
\end{enumerate}
此外, 若 $Z$ 是 $X$ 的子簇, 则限制 $F|_{Z}:Z\to F(Z)$ 是拟射影簇的同构.
\end{Corollary}
\begin{Proof}
Exercise~\ref{ex:isomorphisms-identify-open-closed-and-locally-closed-subsets}.
\end{Proof}

本节余下部分用 Theorem~\ref{thm:regular-maps-are-zariski-continuous} 建立正则映射和同构的其他刻画. 这些刻画只用连续性和正则函数表述, 不涉及拟射影簇所在的环境射影空间. 正如本节最后将讨论的, 这些刻画使我们更接近理解拟射影簇的真正内在性质.

\begin{Proposition}{正则映射的刻画}\leavevmode\label{prop:characterize-regular-maps-by-continuity-and-pullback}
令 $X\subseteq\mathbb{P}^{m}$ 和 $Y\subseteq\mathbb{P}^{n}$ 为拟射影簇. 映射 $F:X\to Y$ 正则, 当且仅当满足以下两个条件:
\begin{enumerate}[label=\arabic*.]
  \item $F$ 关于 Zariski 拓扑连续;
  \item 对每个开集 $V\subseteq Y$ 及每个正则函数 $G\in K[V]$, 函数 $G\circ F$ 在 $F^{-1}(V)$ 上正则.
\end{enumerate}
\end{Proposition}
\begin{Proof}
\BookPairedNote{
先假设映射 $F:X\to Y$ 正则. 则命题陈述中的条件 1 来自 Theorem~\ref{thm:regular-maps-are-zariski-continuous}. 条件 2 来自这样的观察: $F$ 在开集 $F^{-1}(V)$ 上的限制仍正则, 而且正则映射诱导的拉回是良定义的 (Exercise~\ref{ex:pullback-regular-functions-is-algebra-homomorphism}).
}{
见 Exercise~\ref{ex:regular-maps-via-affine-open-covers}, 其中还用仿射开覆盖刻画了正则映射.
}
因此, 只需再证明这一双向陈述的另一个方向.

假设 $F:X\to Y$ 满足命题陈述中的条件 1 和 2. 我们要证明 $F$ 正则. 由 Lemma~\ref{lem:regularity-is-local-on-domain}, 只须证明每个 $p\in X$ 都包含在某个开集 $U_{p}\subseteq X$ 中, 使限制 $F:U_{p}\to Y$ 正则. 令 $p\in X$. 不妨假设 $F(p)\notin\mathcal{V}(y_{0})$, 先定义 $U=F^{-1}(Y\cap\mathbb{A}^{n}_{0})$. 这是 $X$ 的开集 (由条件 1), 且包含 $p$ (因为 $F(p)\notin\mathcal{V}(y_{0})$).

由于对每个 $a\in U$ 有 $F(a)\notin\mathcal{V}(y_{0})$, 可以对每个 $F(a)$ 的齐次坐标作缩放, 使 $y_{0}$ 坐标为 $1$. 因此, 存在唯一的函数 $F_{1},\ldots,F_{n}:U\to K$, 使
\begin{equation}\label{eq:regular-map-in-target-affine-coordinates}
  F(a)=[1:F_{1}(a):\cdots:F_{n}(a)]\qquad\text{对每个 }a\in U.
\end{equation}
换句话说, $F_{i}=(y_{i}/y_{0})\circ F$. 注意 $y_{1}/y_{0},\ldots,y_{n}/y_{0}$ 是 $\mathbb{A}^{n}_{0}$ 上的正则函数. 由条件 2 可知, 各复合 $(y_{i}/y_{0})\circ F=F_{i}$ 在 $U$ 上正则. 由于 $p\in U$, 由正则函数的定义, 对所有 $i=1,\ldots,n$, 都存在同次数的齐次多项式 $g_{i},h_{i}\in K[x_{0},\ldots,x_{m}]$, 使 $h_{i}(p)\ne0$, 且
\begin{equation}\label{eq:target-affine-coordinate-is-local-quotient}
  F_{i}(a)=g_{i}(a)/h_{i}(a)
  \qquad\text{对所有 }a\in U\setminus\mathcal{V}(h_{i}).
\end{equation}
定义 $f_{0}=h_{1}\cdots h_{n}$, 对 $i=1,\ldots,n$ 定义 $f_{i}=f_{0}g_{i}/h_{i}$, 并置 $U_{p}=U\setminus\mathcal{V}(f_{0})$. 则 $U_{p}$ 是包含 $p$ 的 $X$ 的开集. 稍加思考可知多项式 $f_{0},\ldots,f_{n}\in K[x_{0},\ldots,x_{m}]$ 都是同次数的齐次多项式, 而且由 \eqref{eq:regular-map-in-target-affine-coordinates} 和 \eqref{eq:target-affine-coordinate-is-local-quotient} 可知
\[
  F(a)=[f_{0}(a):\cdots:f_{n}(a)]\qquad\text{对所有 }a\in U_{p}.
\]
因此已证明 $F:U_{p}\to Y$ 正则.

这样, 对每个 $p\in X$, 我们都构造了开集 $U_{p}\subseteq X$, 使限制 $F:U_{p}\to Y$ 正则. 因而 $F:X\to Y$ 正则.
\end{Proof}

在上一结果的基础上, 现在用开集及开集上的正则函数刻画拟射影簇之间的同构.

\begin{Corollary}{同构的刻画}\leavevmode\label{cor:characterize-quasiprojective-isomorphisms-by-local-rings-of-functions}
令 $X,Y$ 为拟射影簇. 映射 $F:X\to Y$ 是同构, 当且仅当满足以下两个条件:
\begin{enumerate}[label=\arabic*.]
  \item $F$ 关于 Zariski 拓扑是同胚;
  \item 对每个开集 $V\subseteq Y$, 映射 $G\mapsto G\circ F$ 是从 $K[V]$ 到 $K[F^{-1}(V)]$ 的良定义的 $K$-代数同构.
\end{enumerate}
\end{Corollary}
\begin{Proof}
\BookPairedNote{
先假设 $F$ 为同构. 则 $F$ 是正则映射并有正则逆映射, 由 Theorem~\ref{thm:regular-maps-are-zariski-continuous} 可知 $F$ 是同胚.
}{
见 Exercise~\ref{ex:quasiprojective-isomorphisms-via-affine-open-cover}, 其中还用仿射开覆盖刻画了同构.
}
此外, 对每个开集 $V\subseteq Y$, $F$ 限制为同构 $F:F^{-1}(V)\to V$. 由 Proposition~\ref{prop:quasiprojective-isomorphism-induces-regular-function-ring-isomorphism} 可知 $F^{*}:K[V]\to K[F^{-1}(V)]$ 是同构. 这就证明了双向叙述的一个方向.

反过来, 假设 $F:X\to Y$ 满足命题陈述中的两个条件. 因为 $F$ 连续, 且对每个 $G\in K[V]$ 有 $G\circ F\in K[F^{-1}(V)]$, 由 Proposition~\ref{prop:characterize-regular-maps-by-continuity-and-pullback} 可知 $F$ 正则. 又因为 $F$ 是同胚, 它有连续逆映射 $F^{-1}:Y\to X$. 只需证明 $F^{-1}$ 也正则.

令 $U\subseteq X$ 为开集. 将条件 2 应用于开集 $F(U)\subseteq Y$, 可知 $G\mapsto G\circ F$ 给出同构 $K[F(U)]\to K[U]$. 由此同态的满射性, 对每个 $H\in K[U]$, 都存在 $G\in K[F(U)]$ 使 $H=G\circ F$. 特别地,
\[
  H\circ F^{-1}=G\circ F\circ F^{-1}=G\in K[F(U)].
\]
\BookPairedNoteLeft{
我们的论证没有用到 $G\mapsto G\circ F$ 的单射性. 这一单射性由 $F$ 的满射性推出.
}{
总而言之, 已验证对每个开集 $U\subseteq X$ 和每个正则函数 $H\in K[U]$, 复合 $H\circ F^{-1}$ 在集合 $F(U)=(F^{-1})^{-1}(U)$ 上正则.
}
由于 $F^{-1}$ 也连续, 由 Proposition~\ref{prop:characterize-regular-maps-by-continuity-and-pullback} 可知 $F^{-1}$ 正则, 如所需.
\end{Proof}

到目前为止的发展使我们开始窥见拟射影簇的内在性质. 与仿射簇的情形不同, 仿射簇的内在性质由一个 $K$-代数 (它的坐标环) 完全确定, 而一般拟射影簇的内在性质既有拓扑的一面, 又有代数的一面, 两者都必须考虑. 更具体地说, 拟射影簇 $X$ 的本质数据既包括 $X$ 中开集的信息 ($X$ 上的一个拓扑), 也包括每个开集 $U\subseteq X$ 上正则函数环的信息 ($X$ 上的一个 $K$-代数\textquotedbl{}层\textquotedbl{}). 由 Corollary~\ref{cor:characterize-quasiprojective-isomorphisms-by-local-rings-of-functions}, 拟射影簇之间的函数 $F:X\to Y$ 是同构, 当且仅当 $F$ 是同胚, 即 $F$ 保持拓扑结构, 并且对每个开集 $V\subseteq Y$, $F^{*}:K[V]\to K[F^{-1}(V)]$ 是 $K$-代数的同构, 即 $F$ 保持代数结构.

在仿射情形, 情况简单得多: 要知道仿射情形中的一个映射是否为同构, 只须检查其拉回是否为 $K$-代数的同构, 随即无须另行验证就得到拓扑同胚. 换句话说, 在仿射情形, 代数完全支配拓扑 (以及所有其他内在数据). 然而, 如我们已经看到的 (例如 Example~\ref{ex:regular-functions-on-projective-space-are-constant}), 更一般的拟射影情形并非如此: 整体正则函数的 $K$-代数通常不能刻画簇的全部内在性质. 研究一般的拟射影簇时, 我们需要其拓扑结构, 因为它实质上记录了局部仿射邻域 (在那里, 代数完全支配拓扑) 如何粘合为更整体的对象.

\BookExerciseHeading
\begin{BookExercises}
  \item\label{ex:equivalent-open-and-closed-characterizations-of-continuity} 令 $X,Y$ 为拓扑空间, $F:X\to Y$ 为映射. 证明对 $Y$ 中每个开集 $V$, $F^{-1}(V)$ 在 $X$ 中开, 当且仅当对 $Y$ 中每个闭集 $Z$, $F^{-1}(Z)$ 在 $X$ 中闭.
  \item\label{ex:restrict-domain-or-change-codomain-of-continuous-map} 令 $F:X\to Y$ 为拓扑空间之间的连续映射.
  \begin{enumerate}
    \item 令 $Z$ 为 $X$ 的子空间, 证明限制定义域给出连续映射 $F:Z\to Y$.
    \item 令 $Z$ 为以 $Y$ 为子空间的拓扑空间. 证明扩大陪域给出连续映射 $F:X\to Z$.
    \item 令 $Z$ 为 $Y$ 中包含 $F(X)$ 的子空间. 证明限制陪域给出连续映射 $F:X\to Z$.
  \end{enumerate}
  \item\label{ex:preimage-of-projective-vanishing-set-under-polynomial-map} 令 $X\subseteq\mathbb{P}^{m}$ 为拟射影簇, $f_{0},\ldots,f_{n}\in K[x_{1},\ldots,x_{m}]$ 为同次数的齐次多项式, 使得
  \[
    \mathcal{V}(f_{0},\ldots,f_{n})\cap X=\varnothing.
  \]
  若 $F=[f_{0}:\cdots:f_{n}]:X\to\mathbb{P}^{n}$ 且对某些 $g_{1},\ldots,g_{k}\in K[y_{1},\ldots,y_{n}]$, 有 $Z=\mathcal{V}(g_{1},\ldots,g_{k})\subseteq\mathbb{P}^{n}$, 证明
  \[
    F^{-1}(Z)=X\cap\mathcal{V}\bigl(g_{1}(f_{0},\ldots,f_{n}),\ldots,g_{k}(f_{0},\ldots,f_{n})\bigr).
  \]
  \item\label{ex:isomorphisms-identify-open-closed-and-locally-closed-subsets} 证明 Corollary~\ref{cor:quasiprojective-isomorphisms-identify-subvarieties}.
  \item\label{ex:irreducibility-invariant-under-quasiprojective-isomorphism} 证明不可约性是拟射影簇的内在性质. 换句话说, 若 $X,Y$ 是同构的拟射影簇, 证明 $X$ 不可约当且仅当 $Y$ 不可约.
  \item\label{ex:continuous-image-of-irreducible-subspace-is-irreducible} 令 $F:X\to Y$ 为拓扑空间之间的连续映射. 若 $Z$ 是 $X$ 的不可约子空间, 证明 $F(Z)$ 是 $Y$ 的不可约子空间.
  \item\label{ex:vanishing-locus-of-regular-functions-is-closed} 令 $X$ 为拟射影簇, $S\subseteq K[X]$ 为一族正则函数. 定义
  \[
    \mathcal{V}_{X}(S)=\{a\in X\mid F(a)=0\text{ 对所有 }F\in S\}.
  \]
  证明 $\mathcal{V}_{X}(S)$ 在 $X$ 中闭.
  \item\label{ex:regular-maps-via-affine-open-covers} 令 $F:X\to Y$ 为拟射影簇之间的映射. 证明 $F$ 正则, 当且仅当存在仿射开覆盖
  \[
    X=U_{1}\cup\cdots\cup U_{k}\qquad\text{和}\qquad
    Y=V_{1}\cup\cdots\cup V_{\ell},
  \]
  使得对每个 $i\in\{1,\ldots,k\}$, 都存在 $j\in\{1,\ldots,\ell\}$ 使 $U_{i}\subseteq F^{-1}(V_{j})$, 且限制
  \[
    F|_{U_{i}}:U_{i}\to V_{j}
  \]
  是仿射簇之间的正则映射.
  \item\label{ex:quasiprojective-isomorphisms-via-affine-open-cover} 令 $F:X\to Y$ 为拟射影簇之间的双射. 证明 $F$ 为同构, 当且仅当存在仿射开覆盖
  \[
    X=U_{1}\cup\cdots\cup U_{k},
  \]
  使对每个 $i$, $F(U_{i})$ 在 $Y$ 中开, 且限制 $F$ 给出同构 $U_{i}\cong F(U_{i})$.
\end{BookExercises}
